%% bare_conf.tex
%% V1.4b
%% 2015/08/26
%% by Michael Shell
%% See:
%% http://www.michaelshell.org/
%% for current contact information.
%%
%% This is a skeleton file demonstrating the use of IEEEtran.cls
%% (requires IEEEtran.cls version 1.8b or later) with an IEEE
%% conference paper.
%%
%% Support sites:
%% http://www.michaelshell.org/tex/ieeetran/
%% http://www.ctan.org/pkg/ieeetran
%% and
%% http://www.ieee.org/

%%*************************************************************************
%% Legal Notice:
%% This code is offered as-is without any warranty either expressed or
%% implied; without even the implied warranty of MERCHANTABILITY or
%% FITNESS FOR A PARTICULAR PURPOSE! 
%% User assumes all risk.
%% In no event shall the IEEE or any contributor to this code be liable for
%% any damages or losses, including, but not limited to, incidental,
%% consequential, or any other damages, resulting from the use or misuse
%% of any information contained here.
%%
%% All comments are the opinions of their respective authors and are not
%% necessarily endorsed by the IEEE.
%%
%% This work is distributed under the LaTeX Project Public License (LPPL)
%% ( http://www.latex-project.org/ ) version 1.3, and may be freely used,
%% distributed and modified. A copy of the LPPL, version 1.3, is included
%% in the base LaTeX documentation of all distributions of LaTeX released
%% 2003/12/01 or later.
%% Retain all contribution notices and credits.
%% ** Modified files should be clearly indicated as such, including  **
%% ** renaming them and changing author support contact information. **
%%*************************************************************************


% *** Authors should verify (and, if needed, correct) their LaTeX system  ***
% *** with the testflow diagnostic prior to trusting their LaTeX platform ***
% *** with production work. The IEEE's font choices and paper sizes can   ***
% *** trigger bugs that do not appear when using other class files.       ***                          ***
% The testflow support page is at:
% http://www.michaelshell.org/tex/testflow/



\documentclass[conference]{IEEEtran}
% Some Computer Society conferences also require the compsoc mode option,
% but others use the standard conference format.
%
% If IEEEtran.cls has not been installed into the LaTeX system files,
% manually specify the path to it like:
% \documentclass[conference]{../sty/IEEEtran}
\usepackage{booktabs} 
\usepackage{amsmath} 
\usepackage{amsfonts}
\usepackage{algorithmicx}
\usepackage{algorithm}
\usepackage{algpseudocode}
\usepackage{graphicx}
\usepackage{float}
\usepackage{caption}
\usepackage{subcaption}

\usepackage{amssymb}      % For additional mathematical symbols
\usepackage{algpseudocode}% For pseudocode in algorithms


% Some very useful LaTeX packages include:
% (uncomment the ones you want to load)


% *** MISC UTILITY PACKAGES ***
%
%\usepackage{ifpdf}
% Heiko Oberdiek's ifpdf.sty is very useful if you need conditional
% compilation based on whether the output is pdf or dvi.
% usage:
% \ifpdf
%   % pdf code
% \else
%   % dvi code
% \fi
% The latest version of ifpdf.sty can be obtained from:
% http://www.ctan.org/pkg/ifpdf
% Also, note that IEEEtran.cls V1.7 and later provides a builtin
% \ifCLASSINFOpdf conditional that works the same way.
% When switching from latex to pdflatex and vice-versa, the compiler may
% have to be run twice to clear warning/error messages.






% *** CITATION PACKAGES ***
%
%\usepackage{cite}
% cite.sty was written by Donald Arseneau
% V1.6 and later of IEEEtran pre-defines the format of the cite.sty package
% \cite{} output to follow that of the IEEE. Loading the cite package will
% result in citation numbers being automatically sorted and properly
% "compressed/ranged". e.g., [1], [9], [2], [7], [5], [6] without using
% cite.sty will become [1], [2], [5]--[7], [9] using cite.sty. cite.sty's
% \cite will automatically add leading space, if needed. Use cite.sty's
% noadjust option (cite.sty V3.8 and later) if you want to turn this off
% such as if a citation ever needs to be enclosed in parenthesis.
% cite.sty is already installed on most LaTeX systems. Be sure and use
% version 5.0 (2009-03-20) and later if using hyperref.sty.
% The latest version can be obtained at:
% http://www.ctan.org/pkg/cite
% The documentation is contained in the cite.sty file itself.






% *** GRAPHICS RELATED PACKAGES ***
%
\ifCLASSINFOpdf
  % \usepackage[pdftex]{graphicx}
  % declare the path(s) where your graphic files are
  % \graphicspath{{../pdf/}{../jpeg/}}
  % and their extensions so you won't have to specify these with
  % every instance of \includegraphics
  % \DeclareGraphicsExtensions{.pdf,.jpeg,.png}
\else
  % or other class option (dvipsone, dvipdf, if not using dvips). graphicx
  % will default to the driver specified in the system graphics.cfg if no
  % driver is specified.
  % \usepackage[dvips]{graphicx}
  % declare the path(s) where your graphic files are
  % \graphicspath{{../eps/}}
  % and their extensions so you won't have to specify these with
  % every instance of \includegraphics
  % \DeclareGraphicsExtensions{.eps}
\fi
% graphicx was written by David Carlisle and Sebastian Rahtz. It is


% correct bad hyphenation here
\hyphenation{op-tical net-works semi-conduc-tor}


\begin{document}
%
% paper title
% Titles are generally capitalized except for words such as a, an, and, as,
% at, but, by, for, in, nor, of, on, or, the, to and up, which are usually
% not capitalized unless they are the first or last word of the title.
% Linebreaks \\ can be used within to get better formatting as desired.
% Do not put math or special symbols in the title.
\title{Explainable Disease Progression and Counterfactual Video Generation with Vision–Language Models}


% author names and affiliations
% use a multiple column layout for up to three different
% affiliations
%\author{\IEEEauthorblockN{Belkacem Chikhaoui}
%\IEEEauthorblockA{Applied Artificial Intelligence %Institute\\
%Department of Science and Technology\\
%TELUQ University, Montreal, Canada\\
%Email: belkacem.chikhaoui@teluq.ca}

\author{\IEEEauthorblockN{Mohamed Boufafa, Belkacem Chikhaoui, Belkacem Khaldi}
\IEEEauthorblockA{Applied Artificial Intelligence Institute\\
TELUQ University, Montreal, Canada \\
belkacem.chikhaoui@teluq.ca
}

}


% make the title area
\maketitle

% As a general rule, do not put math, special symbols or citations
% in the abstract
\begin{abstract}


\end{abstract}

\IEEEpeerreviewmaketitle

\section{Introduction}


\section{Related Work}




\section{Methodology}
\label{sec:methodology}


\begin{figure*}[h]
\includegraphics[width=1\textwidth]{figure/qtlstm_circuit1.png}
\caption{QLSTM quantum circuit example.} \label{qlstm}
\end{figure*}



\section{Experiments}
\label{sec:experiments}

\subsection{Datasets}

We evaluated the QLSTM on the following datasets:

\begin{itemize}
    \item \textbf{Synthetic Sine Waves:} A collection of artificially generated periodic signals used to assess the model's ability to capture recurring temporal dynamics. Each sequence is generated with varying amplitudes, frequencies, and phase shifts to simulate diverse periodic behaviors. These signals are noise-free or optionally perturbed with Gaussian noise. The dataset is generated programmatically using NumPy.

    \item \textbf{Airline Passenger Dataset:} A univariate time series dataset containing monthly totals of international airline passengers (in thousands) from 1949 to 1960. This dataset is commonly used to evaluate models on trend and seasonality detection. It features strong yearly cycles and a non-linear upward trend. It is available in the \texttt{statsmodels} and \texttt{seaborn} Python libraries and can also be downloaded from Kaggle or directly from the UCI Machine Learning Repository.

    \item \textbf{Daily Minimum Temperatures in Melbourne:} This dataset records the daily minimum temperatures (in degrees Celsius) in Melbourne, Australia, from January 1, 1981, to December 31, 1990. It consists of 3650 time steps and is ideal for evaluating models on long, continuous, and real-world climate-related sequences. It is publicly available from the UCI Machine Learning Repository and is often included in the \texttt{pandas} and \texttt{sklearn.datasets} libraries for time series experiments.

    \item \textbf{COVID-19 Daily Cases:} This dataset contains the daily number of confirmed COVID-19 cases reported in Canada from January 2020 to December 2022. The data includes fields such as date, province, number of daily cases, and cumulative totals. It is used to test how well models capture abrupt changes and long-range dependencies in epidemiological data. The dataset is downloaded through the Johns Hopkins University COVID-19 GitHub repository.

    \item \textbf{Exchange Rate Dataset:} A multivariate time series dataset containing daily exchange rates between various currency pairs, such as USD/GBP, USD/EUR, and USD/JPY, typically from 2000 to 2022. This dataset is highly non-stationary and exhibits volatility clustering, making it ideal for stress-testing time series models. The data is publicly accessible. The dataset is downloaded through Yahoo Finance API.
    
    \item \textbf{Apple Stock Price Dataset:} A univariate time series dataset containing daily closing prices of Apple Inc. (AAPL) stock, sourced from Yahoo Finance. This dataset spans from January 2018 to December 2023 and reflects real-world financial market behavior, including noise, volatility, and non-stationary patterns. It is commonly used to evaluate models on forecasting accuracy under real market conditions. The dataset exhibits short- and long-term trends, occasional abrupt changes (e.g., due to earnings reports or market shocks), and no fixed seasonality. This type of financial data is particularly suitable for testing the effectiveness of quantum-enhanced models in capturing complex temporal dependencies. The dataset is downloaded through Yahoo Finance API.
\end{itemize}


\subsection{Implementation}
The QLSTM is implemented using \texttt{PennyLane} and \texttt{PyTorch}, utilizing 7 qubits and 2 layers of parameterized rotations with entanglement via chained CNOT gates in each quantum circuit. The classical LSTM baseline is implemented with a comparable architecture using 32 hidden units. Both models are trained under the same configuration and evaluated using the Root Mean Squared Error (RMSE) metric to ensure a fair comparison.

\subsection{Results}


\section{Discussion}
\label{sec:discussion}



\section{Conclusion}
\label{sec:conclusion}




\bibliographystyle{IEEEtran}
\bibliography{ref}
\end{document}


